%!TEX root=../../note.tex
\subsection{Proving Uniqueness}
To prove that there is a unique $x\in S$ satisfying $P(x)$, show both that
such an $x$ exists and that any two elements satisfying $P$ are equal.
Proving only the second part does not establish existence.

\begin{proposition}
    For all $a,b\in\Z$, if $a\neq 0$ and $a\mid b$, then there is a unique
    integer $k$ such that $b=ka$.
\end{proposition}
\begin{remark}[Proof idea]
    Divisibility gives existence directly. For uniqueness, compare two
    possible multiples of the same nonzero integer $a$.
\end{remark}
\begin{proof}
    Let $a,b\in\Z$ be arbitrary. Assume $a\neq 0$ and $a\mid b$.

    \textbf{Existence.} By the definition of divisibility, there is some
    $k\in\Z$ such that $b=ka$.

    \textbf{Uniqueness.} Suppose $k,m\in\Z$ both satisfy the equation. Then
    $ka=b=ma$. Since $a\neq 0$, dividing by $a$ gives $k=m$.
\end{proof}
\begin{remark}
    The condition $a\neq 0$ is needed for uniqueness. If $a=b=0$, then
    $0\mid 0$, but every integer $k$ satisfies $0=k\cdot0$.
\end{remark}

\subsubsection*{In practice}
\begin{itemize}
    \item \textbf{Existence, then uniqueness.} First produce (or cite) one
    object with the property. Then assume two objects $x$ and $y$ both have it
    and show $x=y$. Alternatively, assume two \emph{different} objects have it
    and derive a contradiction.
\end{itemize}

\subsubsection*{Exercises}
\begin{enumerate}
    \item Use a proof by contradiction to show that two non-parallel lines
    have a unique intersection point.
\end{enumerate}
